% Auto-generated by code/3_main_analysis/3_appendix/7_semantic_gap_text_interpretability.py -- do not edit manually
\begin{tabular}{lcp{0.26\textwidth}p{0.26\textwidth}p{0.25\textwidth}}
\toprule
SDG & Gap & Research-side distinctive terms & Policy-side distinctive terms & Descriptive reading \\
\midrule
SDG 17 & 0.659 & cells, model, learning, brain, expression, neural, disease, protein & countries, national, country, nations, united, global, agenda, government & High-gap case: the research-side sample is technical and biomedical/ML-heavy, while policy language is international, institutional, and agenda-oriented. \\
SDG 13 & 0.536 & model, learning, language, neural, method, machine, prediction, performance & climate, change, countries, emissions, national, global, nations, action & High-gap case: both sides concern climate, but policy language is more institutional and commitment-oriented while research language is more technical and measurement-oriented. \\
SDG 9 & 0.220 & learning, model, machine, analysis, neural, network, detection, method & national, government, public, infrastructure, countries, strategy, sector, country & Lower-gap comparison: differences remain visible, but the contrast is closer to a technology/infrastructure register split than to the broader institutional gap seen in SDGs 13 and 17. \\
\bottomrule
\end{tabular}
